% Generated by book/tools/notebook_to_tex.py. Do not edit by hand.

\chapter[토크나이저와 데이터셋 (Tokenizer \& Datasets)]{토크나이저와 데이터셋 (Tokenizer \& Datasets)}

\markboth{토크나이저와 데이터셋 (Tokenizer \& Datasets)}{토크나이저와 데이터셋 (Tokenizer \& Datasets)}

\chaptermeta{08_tokenizer_datasets/08_tokenizer_datasets.ipynb}{https://colab.research.google.com/github/yoon-gu/neuqes-101/blob/master/08_tokenizer_datasets/08_tokenizer_datasets.ipynb}{padding, truncation, max\_length와 datasets 입력 파이프라인}



\index{1Tokenizer@Tokenizer}

\index{1datasets@datasets}

\index{1load dataset@load\_dataset}

\index{1Dataset map@Dataset.map}

\index{1Dataset filter@Dataset.filter}

\index{1padding@padding}

\index{1truncation@truncation}

\index{1max length@max\_length}

\index{1attention mask@attention\_mask}

\index{1DataCollatorWithPadding@DataCollatorWithPadding}

\index{1DataLoader@DataLoader}

\index{1Apache Arrow@Apache Arrow}

\index{0토크나이저 옵션@토크나이저 옵션}

\index{0패딩@패딩}

\index{0잘림@잘림}

\index{0최대 길이@최대 길이}

\index{0어텐션 마스크@어텐션 마스크}

\index{0데이터셋@데이터셋}

\index{0입력 파이프라인@입력 파이프라인}

\index{1memory mapping@memory mapping}

\index{1Dataset select@Dataset.select}

\index{1Dataset shuffle@Dataset.shuffle}

\index{1with format@with\_format}

\index{1torch format@torch format}

\index{1token length distribution@token length distribution}

\index{195th percentile@95th percentile}

\index{1padding True@padding=True}

\index{1padding max length@padding="max\_length"}

\index{1truncation True@truncation=True}

\index{1input ids@input\_ids}

\index{0토큰 길이 분포@토큰 길이 분포}

\index{0메모리 매핑@메모리 매핑}

\index{0배치 패딩@배치 패딩}

\index{0고정 길이 패딩@고정 길이 패딩}

\index{0동적 패딩@동적 패딩}



\section{변화추적표}

\begin{booktable}{8장 변화추적표}{tab:ch08-01}
\begin{adjustbox}{max width=\textwidth}
\begin{tabular}{@{}lllllll@{}}
\toprule
Ch & 모델 & 토크나이저 & 데이터 & Output Head & Activation & Loss \\
\midrule

1 & (TF-IDF) & \inlinecode{TfidfVectorizer()} & Yelp 5,000 & --- & --- & --- \\
2-6 & sklearn 모델들 & \inlinecode{TfidfVectorizer()} & Yelp 변형 & 1차원/K차원 & 없음/sigmoid/softmax & MSE/BCE/CE \\
7 & \inlinecode{pipeline("sentiment-analysis")} & \inlinecode{AutoTokenizer.from\_pretrained(...)} & 간단 영어 예시 & 사전학습 헤드 & softmax & --- \\
\textbf{8 ← 여기} & (모델 없음 --- 토크나이저·데이터 파이프라인만) & \inlinecode{AutoTokenizer.from\_pretrained(...)} & \textbf{Yelp 5,000 (Phase 0과 동일)} & --- & --- & --- \\
\bottomrule
\end{tabular}
\end{adjustbox}
\end{booktable}

전체 20장 표는 \href{https://github.com/yoon-gu/neuqes-101\#챕터별-변화추적표}{루트 README.md}를 참고하세요.



\section{변경점: 7장 대비}

\begin{booktable}{8장 변경점 요약}{tab:ch08-02}
\begin{adjustbox}{max width=\textwidth}
\begin{tabular}{@{}lll@{}}
\toprule
축 & 7장 & 8장 \\
\midrule

모델 & \inlinecode{pipeline} + \inlinecode{AutoModelForSequenceClassification} & \textbf{모델 로드 없음} (다음 장 학습 준비 단계) \\
토크나이저 & WordPiece --- 한 문장 시연 & \textbf{WordPiece + 옵션 학습} (\inlinecode{padding} / \inlinecode{truncation} / \inlinecode{max\_length}) \\
데이터 & 간단 영어 예시 문장 & \textbf{\inlinecode{datasets} 로 Yelp 65만 → 5,000 subsample} (Phase 0과 동일 데이터) \\
데이터 라이브러리 & (없음) & \textbf{\inlinecode{datasets}} 첫 등장 --- \inlinecode{load\_dataset}, \inlinecode{map}, \inlinecode{filter}, \inlinecode{with\_format} \\
학습 단계 & 추론만 & 학습·추론 모두 없음 --- 데이터 파이프라인 \emph{연습} \\
\bottomrule
\end{tabular}
\end{adjustbox}
\end{booktable}

\textbf{왜 이 장?} 9장 BERT 회귀에서 \inlinecode{Trainer.train()} 한 줄을 부르려면, 그 한 줄에 어떤 입력이 들어가는지 미리 알고 있어야 합니다. \inlinecode{Dataset} 객체, \inlinecode{padding}/\inlinecode{truncation} 결정, \inlinecode{DataLoader} 변환이 그 한 줄을 떠받치는 부품들입니다. 이 장은 그 입력 형태를 \emph{학습 없이} 미리 손에 익히는 자리입니다.

\textbf{Phase 0와의 다리}: Yelp 5,000건은 1-6장에서 줄곧 쓴 데이터. 같은 텍스트가 TF-IDF에서 sparse vector로 갔던 길이, 이번엔 WordPiece에서 \inlinecode{input\_ids} + \inlinecode{attention\_mask} 텐서 쌍으로 가는 길을 봅니다.



\section{\texorpdfstring{토크나이저 노트}{토크나이저 노트}}

WordPiece가 출력하는 시퀀스 길이는 입력 텍스트마다 다릅니다. 그런데 모델은 \textbf{고정된 shape의 텐서 배치} 를 받아야 하므로, 이 둘 사이를 맞추는 세 가지 옵션이 있습니다.

\begin{booktable}{8장 토크나이저 노트 표}{tab:ch08-03}
\begin{adjustbox}{max width=\textwidth}
\begin{tabular}{@{}lll@{}}
\toprule
옵션 & 의미 & 언제 쓰나 \\
\midrule

\inlinecode{padding=False} & 패딩 없음 (기본값). 시퀀스 길이가 다 다름 & 한 문장씩 처리할 때 \\
\inlinecode{padding=True} & \textbf{배치 안 가장 긴 시퀀스 길이까지} padding & 일반 학습 (효율적, dynamic padding) \\
\inlinecode{padding="max\_length"} & \textbf{항상 \inlinecode{max\_length} 까지} padding (짧으면 패딩, 길면 자름) & TPU·고정 shape 필요할 때 \\
\inlinecode{truncation=True} & \inlinecode{max\_length} 초과분은 \emph{잘라냄} & 항상 같이 두는 게 안전 (긴 입력 방지) \\
\inlinecode{max\_length=N} & 길이 상한 (모델별 사전학습 한도 --- BERT 512) & 메모리/속도 trade-off \\
\bottomrule
\end{tabular}
\end{adjustbox}
\end{booktable}

\textbf{패딩이 들어간 자리는 \inlinecode{attention\_mask=0}} 으로 표시됩니다. 모델은 이 mask를 보고 self-attention에서 패딩 토큰을 무시하므로, 아무리 길게 패딩을 붙여도 학습 결과는 달라지지 않습니다. 다만 그만큼 속도와 메모리만 낭비될 뿐입니다.

이 장에서 위 세 옵션을 직접 호출해 input\_ids와 attention\_mask가 어떻게 변하는지 봅니다.



\Needspace{8\baselineskip}
\begin{lstlisting}
!pip install -q transformers datasets
\end{lstlisting}

\begin{codeRead}
\textbf{1행}의 \inlinecode{!pip install -q transform...}에서는 Colab 실행에 필요한 패키지를 설치합니다.
\end{codeRead}



\Needspace{24\baselineskip}
\begin{lstlisting}
import warnings
warnings.filterwarnings("ignore")

import numpy as np
import pandas as pd
import torch
from datasets import load_dataset
from transformers import AutoTokenizer

print(f"PyTorch:        {torch.__version__}")
print(f"CUDA available: {torch.cuda.is_available()}")
if torch.cuda.is_available():
    print(
        f"GPU:            "
        f"{torch.cuda.get_device_name(0)}"
    )
print(
    "\nNo model weights loaded in this chapter; VRAM "
    "stays roughly flat."
)
\end{lstlisting}

\begin{codeRead}
\inlinecode{numpy, pandas, PyTorch, datasets, transformers} 같은 주요 패키지를 불러옵니다.
\end{codeRead}



\section{\texorpdfstring{데이터셋 로드}{데이터셋 로드}}

\inlinecode{load\_dataset("yelp\_review\_full")} 한 줄로 Hugging Face Hub에서 65만 건 학습 데이터를 받아옵니다 (50K test). 처음 받으면 \textasciitilde150MB 다운로드 + 디스크 캐시.

\textbf{주목할 점}: \inlinecode{datasets} 는 Apache Arrow 형식으로 디스크에 저장하고 메모리맵으로 접근합니다. 65만 건이 한꺼번에 RAM에 올라가는 게 아니라, 인덱싱하는 시점에만 디스크에서 필요한 부분을 읽어 옵니다. 그래서 데이터셋이 아무리 커도 RAM 사용량에는 거의 영향이 없습니다.



\Needspace{8\baselineskip}
\begin{lstlisting}
ds = load_dataset("yelp_review_full")
print(ds)
\end{lstlisting}



\Needspace{13\baselineskip}
\begin{lstlisting}
print(f"train samples: {len(ds['train']):,}")
print(f"test samples:  {len(ds['test']):,}")
print(f"\nfeatures: {ds['train'].features}")
print(f"\nFirst sample:")
print(
    f"  label: "
    f"{ds['train'][0]['label']}  (0-4 = stars 1-5)"
)
print(f"  text:  {ds['train'][0]['text'][:200]}...")
\end{lstlisting}



\Needspace{8\baselineskip}
\begin{lstlisting}
small = ds["train"].shuffle(seed=42).select(range(5000))
print(small)
print(f"\nfirst sample text: {small[0]['text'][:150]}...")
\end{lstlisting}



\section{토크나이저 옵션}

7장과 같은 \inlinecode{distilbert-base-uncased} 토크나이저로 시작합니다 (사전학습 모델 그대로).



\Needspace{11\baselineskip}
\begin{lstlisting}
tokenizer = AutoTokenizer.from_pretrained("distilbert-base-uncased")
print(f"Class:     {type(tokenizer).__name__}")
print(f"vocab:     {tokenizer.vocab_size:,}")
print(
    f"pad_token: {tokenizer.pad_token}  (id="
    f"{tokenizer.pad_token_id})"
)
\end{lstlisting}

\begin{codeRead}
\textbf{1행}의 \inlinecode{tokenizer = ...}에서는 이후 분석에서 사용할 값을 준비합니다.
\end{codeRead}



\subsection{기본 토큰화}

\Needspace{6\baselineskip}
\begin{lstlisting}
tokenizer(text)
\end{lstlisting}



\Needspace{14\baselineskip}
\begin{lstlisting}
sample = small[0]["text"]
print(f"Input (first 150 chars): {sample[:150]}...\n")

out = tokenizer(sample)
print(f"input_ids length: {len(out['input_ids'])}")
print(f"First 30 IDs:      {out['input_ids'][:30]}")
print(
    f"Decoded first 30:  "
    f"{tokenizer.decode(out['input_ids'][:30])}"
)
\end{lstlisting}



\subsection{\texorpdfstring{동적 패딩}{동적 패딩}}

여러 문장을 한 배치로 묶으려면 길이가 모두 같아야 텐서가 만들어집니다. \inlinecode{padding=True} 옵션을 주면 한 배치 안에서 \textbf{가장 긴 문장 길이에 맞춰 짧은 문장만 패딩으로 채우므로}, 배치마다 필요한 만큼만 늘어나 가장 효율적입니다.



\Needspace{26\baselineskip}
\begin{lstlisting}
short_text = "Great service!"
long_text = small[0]["text"]
texts = [short_text, long_text]

out_no_pad = tokenizer(texts, padding=False)
print(f"padding=False:")
for i, ids in enumerate(out_no_pad["input_ids"]):
    print(f"  sentence {i}: {len(ids)} tokens")

out_dyn = tokenizer(
    texts,
    padding=True,
    return_tensors="pt",
)
print(f"\npadding=True (return_tensors='pt'):")
print(f"  input_ids shape: {out_dyn['input_ids'].shape}")
print(
    f"  attention_mask sentence 0: "
    f"{out_dyn['attention_mask'][0][:20]}"
)
print(
    f"  attention_mask sentence 1: "
    f"{out_dyn['attention_mask'][1][:20]}"
)
\end{lstlisting}



\subsection{\texorpdfstring{고정 길이 패딩}{padding=“max\_length”, max\_length=128 --- 고정 길이}}

배치마다 길이가 달라지는 게 싫을 때 (TPU·정적 그래프 환경) 항상 \inlinecode{max\_length} 까지 padding합니다.



\Needspace{21\baselineskip}
\begin{lstlisting}
out_fixed = tokenizer(
    texts,
    padding="max_length",
    max_length=128,
    return_tensors="pt",
)
print(
    f"shape: "
    f"{out_fixed['input_ids'].shape}  (batch 2, max_length=128)"
)

real_ratio = out_fixed["attention_mask"].sum().item() / out_fixed["attention_mask"].numel()
print(f"\nattention_mask=1 ratio: {real_ratio:.1%}")
print(
    f"  → short sentence is mostly padding (compute "
    f"wasted)"
)
\end{lstlisting}



\subsection{\texorpdfstring{긴 입력 자르기}{긴 입력 자르기}}

BERT 계열은 사전학습 시 \inlinecode{max\_length=512} 로 학습돼서 그보다 긴 입력은 처리할 수 없습니다. \inlinecode{truncation=True} 로 자동 절단.



\Needspace{26\baselineskip}
\begin{lstlisting}
very_long = "Hello world! This is a sentence. " * 200

out_full = tokenizer(very_long)
print(
    f"truncation=False: "
    f"{len(out_full['input_ids'])} tokens  (may exceed BERT limit 512)"
)

# truncation=True + max_length=128
out_trunc = tokenizer(
    very_long,
    truncation=True,
    max_length=128,
)
print(
    f"truncation=True, max_length=128: "
    f"{len(out_trunc['input_ids'])} tokens"
)
print(
    f"  Last token: "
    f"{tokenizer.decode([out_trunc['input_ids'][-1]])} (= [SEP], always appended)"
)
\end{lstlisting}



\subsection{attention mask}

핵심: 패딩 토큰이 \emph{다른 토큰의 표현에 영향을 주지 않도록} 막습니다.

\Needspace{8\baselineskip}
\begin{lstlisting}
attention_scores = Q @ K.T / sqrt(d_k)
attention_scores[mask == 0] = -inf
attention_weights = softmax(attention_scores)
output = attention_weights @ V
\end{lstlisting}

\inlinecode{-inf} 가 softmax를 거치면 \texttt{e\^{}(-inf)\ =\ 0} 이 되어 패딩 토큰의 가중치가 정확히 0이 됩니다. 결과적으로 아무리 길게 패딩을 붙여도 학습 결과는 변하지 않고, 그만큼의 계산이 그저 버려질 뿐입니다.

확인: 위 padding=“max\_length” 출력에서 attention\_mask=1 비율이 낮으면 그만큼 \emph{낭비된 계산} 입니다.



\subsection{\texorpdfstring{max\_length 결정}{max\_length 결정}}

너무 작으면 \emph{정보 손실} (긴 리뷰가 잘림), 너무 크면 \emph{낭비} (대부분 패딩). 실제 데이터의 토큰 길이 분포를 보고 정합니다.



\Needspace{26\baselineskip}
\begin{lstlisting}
lengths = []
for i in range(min(1000, len(small))):
    n = len(tokenizer.tokenize(small[i]["text"]))
    lengths.append(n)
lengths = np.array(lengths)

print(f"Token length distribution over 1,000 samples:")
print(f"  min:    {lengths.min()}")
print(f"  mean:   {lengths.mean():.0f}")
print(f"  median: {int(np.median(lengths))}")
print(f"  p90:    {int(np.percentile(lengths, 90))}")
print(f"  p95:    {int(np.percentile(lengths, 95))}")
print(f"  p99:    {int(np.percentile(lengths, 99))}")
print(f"  max:    {lengths.max()}")

print(f"\nFraction truncated at various max_length:")
for max_len in [64, 128, 256, 512]:
    truncated_pct = (lengths > max_len).mean() * 100
    print(
        f"  max_length={max_len}: "
        f"{truncated_pct:5.1f}% truncated"
    )
\end{lstlisting}



\textbf{해석}: 학습 시 \inlinecode{max\_length=128} 로 두면 95\% 이상 정상, 5\% 정도만 잘립니다 (Yelp 리뷰가 대부분 짧음). \inlinecode{max\_length=512} 면 거의 모든 리뷰를 보존하지만 평균 패딩 비율이 60-70\%라 메모리·시간 낭비.

이 커리큘럼은 \textbf{\inlinecode{max\_length=128}} 을 표준으로 씁니다 (T4 30분 제약 + 무난한 정보 보존 균형).



\section{\texorpdfstring{일괄 토큰화}{일괄 토큰화}}

샘플 하나씩 \inlinecode{tokenizer(...)} 부르는 건 5,000번 호출이 필요하고 느립니다. \inlinecode{dataset.map(fn,\ batched=True)} 로 \emph{배치 단위} 일괄 호출이 표준.



\Needspace{26\baselineskip}
\begin{lstlisting}
def tokenize_fn(batch):

    return tokenizer(
        batch["text"],
        padding="max_length",
        truncation=True,
        max_length=128,
    )

tokenized = small.map(
    tokenize_fn,
    batched=True,
    batch_size=200,
)

print(tokenized)
print(
    f"\nFirst sample input_ids length: "
    f"{len(tokenized[0]['input_ids'])}  (= 128, fixed)"
)
print(
    f"First sample attention_mask sum: "
    f"{sum(tokenized[0]['attention_mask'])}  (real tokens)"
)
\end{lstlisting}



\subsection{\texorpdfstring{샘플 필터링}{샘플 필터링}}



\Needspace{15\baselineskip}
\begin{lstlisting}
positive = small.filter(lambda x: x["label"] >= 3)
print(
    f"Positive samples: {len(positive):,} / {len(small):,} = "
    f"{len(positive)/len(small):.1%}"
)

short = small.filter(lambda x: len(x["text"].split()) <= 100)
print(
    f"Short samples:    {len(short):,} / {len(small):,} = "
    f"{len(short)/len(small):.1%}"
)
\end{lstlisting}



\subsection{\texorpdfstring{텐서 형식 변환}{with\_format(“torch”) --- 텐서 형식으로}}



\Needspace{22\baselineskip}
\begin{lstlisting}
tokenized_torch = tokenized.with_format(
    "torch",
    columns=["input_ids", "attention_mask", "label"],
)

sample = tokenized_torch[0]
print(
    f"input_ids:      {type(sample['input_ids']).__name__}, dtype={sample['input_ids'].dtype}, shape="
    f"{sample['input_ids'].shape}"
)
print(
    f"attention_mask: {type(sample['attention_mask']).__name__}, shape="
    f"{sample['attention_mask'].shape}"
)
print(
    f"label:          "
    f"{sample['label']}  (0-4 = stars 1-5)"
)
\end{lstlisting}



\section{\texorpdfstring{DataLoader 변환}{DataLoader 변환}}

PyTorch \inlinecode{DataLoader} 는 dataset을 받아 \emph{배치 + shuffle} 을 자동 처리합니다. 9장의 \inlinecode{Trainer} 가 내부에서 이걸 만들어 쓰지만, 직접 만들 줄 알면 디버깅에 유리.



\Needspace{24\baselineskip}
\begin{lstlisting}
from torch.utils.data import DataLoader

loader = DataLoader(
    tokenized_torch,
    batch_size=8,
    shuffle=True,
)

batch = next(iter(loader))
print(f"batch keys:           {list(batch.keys())}")
print(
    f"input_ids shape:      "
    f"{batch['input_ids'].shape}      (= [batch_size, max_length])"
)
print(
    f"attention_mask shape: "
    f"{batch['attention_mask'].shape}"
)
print(f"label shape:          {batch['label'].shape}")
print(f"label values:         {batch['label'].tolist()}")
\end{lstlisting}



\section{\texorpdfstring{DataCollator}{DataCollator}}

위 \inlinecode{DataLoader} 코드는 \inlinecode{padding="max\_length"} 로 \emph{모든} 샘플을 128로 미리 padding한 상태였습니다. 짧은 문장은 대부분 패딩 자리라 메모리 낭비가 큽니다.

\textbf{더 좋은 방법: 토큰화 시엔 padding 안 하고}, \inlinecode{DataLoader} 가 \emph{배치를 만들 때마다 그 배치 안 longest까지만} padding (동적 padding). 이걸 담당하는 게 \inlinecode{DataCollator}.

\inlinecode{DataCollator} 는 \inlinecode{DataLoader} 의 \inlinecode{collate\_fn} 자리에 들어가는 함수입니다. 매 배치마다 N개 샘플을 받아 batch-level 변환을 적용한 뒤 텐서로 묶어주는 역할을 하고, Hugging Face는 task별로 여러 종류를 제공합니다.



\Needspace{22\baselineskip}
\begin{lstlisting}
from transformers import DataCollatorWithPadding

def tokenize_dynamic(batch):
    return tokenizer(
        batch["text"],
        truncation=True,
        max_length=128,
    )

tokenized_dyn = small.select(range(50)).map(tokenize_dynamic, batched=True)
tokenized_dyn = tokenized_dyn.remove_columns(["text"])

sample_lens = [len(tokenized_dyn[i]["input_ids"]) for i in range(10)]
print(f"First 10 sample token lengths: {sample_lens}")
print(
    f"  → all different — cannot batch as-is into a "
    f"tensor"
)
\end{lstlisting}



\Needspace{26\baselineskip}
\begin{lstlisting}
collator = DataCollatorWithPadding(
    tokenizer=tokenizer,
    padding=True,
)
dyn_loader = DataLoader(
    tokenized_dyn, batch_size=8, shuffle=False,
    collate_fn=collator,
)

print(f"Per-batch shape (batch_size=8, varies):")
print(
    f"{'batch':>6}  {'shape':>16}  {'real_tok':>10}  {'total':>6}  "
    f"{'fill':>6}"
)
for i, batch in enumerate(dyn_loader):
    real = batch["attention_mask"].sum().item()
    total = batch["attention_mask"].numel()
    print(
        f"{i:>6}  {str(tuple(batch['input_ids'].shape)):>16}  {real:>10}  {total:>6}  "
        f"{real/total:>6.0%}"
    )
    if i >= 4: break
\end{lstlisting}



\textbf{정적 vs 동적 padding 비교}

\begin{booktable}{8장 DataCollator 표}{tab:ch08-04}
\begin{adjustbox}{max width=\textwidth}
\begin{tabular}{@{}lllll@{}}
\toprule
방식 & 토큰화 시점 & 배치 shape & 실제 토큰 비율 & 메모리/속도 \\
\midrule

\textbf{정적} (\inlinecode{padding="max\_length"=128}) & 모든 샘플을 128로 미리 padding & 항상 (B, 128) & 보통 30-50\% & 일정, 낭비 큼 \\
\textbf{동적} (\inlinecode{DataCollatorWithPadding}) & padding 없이 truncation만 & (B, longest in batch) --- 매번 다름 & 보통 70-95\% & 변동, 효율 \\
\bottomrule
\end{tabular}
\end{adjustbox}
\end{booktable}

위 출력에서 채움률(실제 토큰/전체)이 정적 방식보다 훨씬 높을 겁니다. 같은 학습이라도 동적 padding이 \emph{2배 가까이 빠른 경우} 도 흔합니다.

\subsection{\texorpdfstring{다양한 \inlinecode{DataCollator} 종류}{다양한 DataCollator 종류}}

\begin{booktable}{8장 다양한 DataCollator 종류 표}{tab:ch08-05}
\begin{adjustbox}{max width=\textwidth}
\begin{tabular}{@{}lll@{}}
\toprule
Collator & 용도 & 자주 쓰이는 곳 \\
\midrule

\inlinecode{DataCollatorWithPadding} & 분류·회귀 --- \inlinecode{input\_ids}/\inlinecode{attention\_mask} 동적 padding & \textbf{9-13장 모든 분류 학습 (기본)} \\
\inlinecode{DataCollatorForLanguageModeling} & MLM --- 입력의 15\%를 \texttt{{[}MASK{]}} 로 가려 라벨 생성 & BERT 사전학습 재현, MLM 헤드 학습 \\
\inlinecode{DataCollatorForSeq2Seq} & seq2seq --- encoder/decoder input 둘 다 padding & T5, BART 같은 인코더-디코더 학습 \\
\inlinecode{DataCollatorForTokenClassification} & NER 같은 토큰 단위 라벨링 --- labels도 padding & NER, POS tagging \\
\inlinecode{default\_data\_collator} & 단순 stacking --- padding 없이 길이 같은 샘플들에 & 이미 padding 끝낸 데이터 \\
\bottomrule
\end{tabular}
\end{adjustbox}
\end{booktable}

\subsection{학습 코드와의 연결}

이 장에서 만든 \emph{데이터 파이프라인 부품들}이 다음 장부터 학습 코드에서 어떤 자리에 들어가는지 미리 그려두면, 9장 이후 코드가 훨씬 익숙해 보입니다.

\subsubsection{\texorpdfstring{패턴 A --- \inlinecode{Trainer} (커리큘럼 기본, 9-13장 대부분)}{패턴 A --- Trainer (커리큘럼 기본, 9-13장 대부분)}}

\Needspace{24\baselineskip}
\begin{lstlisting}
tokenizer = AutoTokenizer.from_pretrained("distilbert-base-uncased")
def tok(b): return tokenizer(b["text"], truncation=True, max_length=128)
train_ds = small.map(
    tok,
    batched=True).remove_columns(["text"],
)

model = AutoModelForSequenceClassification.from_pretrained(
    "distilbert-base-uncased", num_labels=1, problem_type="regression",
)

trainer = Trainer(
    model=model,
    args=TrainingArguments(...),
    train_dataset=train_ds,
    eval_dataset=eval_ds,
    tokenizer=tokenizer,

)
trainer.train()
\end{lstlisting}

\textbf{\inlinecode{tokenizer=...} 한 줄이 collator를 자동 생성합니다.} 학습자가 명시적으로 \inlinecode{data\_collator=DataCollatorWithPadding(...)} 를 적을 일이 거의 없는 이유.

\subsubsection{패턴 B --- 직접 학습 루프 (커스텀이 필요할 때)}

\Needspace{24\baselineskip}
\begin{lstlisting}
train_ds = small.map(
    tok,
    batched=True).remove_columns(["text"],
)

collator = DataCollatorWithPadding(
    tokenizer=tokenizer,
    padding=True,
)
loader = DataLoader(
    train_ds,
    batch_size=16,
    shuffle=True,
    collate_fn=collator,
)

for batch in loader:
    batch = {k: v.to(model.device) for k, v in batch.items()}
    outputs = model(**batch)
    loss = outputs.loss
    loss.backward()
    optimizer.step()
    optimizer.zero_grad()
\end{lstlisting}

\subsubsection{두 패턴의 매핑}

\begin{booktable}{8장 학습 코드와의 연결 표}{tab:ch08-06}
\begin{adjustbox}{max width=\textwidth}
\begin{tabular}{@{}lll@{}}
\toprule
이번 8장에서 만든 것 & 패턴 A (Trainer) & 패턴 B (직접 루프) \\
\midrule

\inlinecode{tokenizer\ =\ AutoTokenizer.from\_pretrained(...)} & \inlinecode{Trainer(tokenizer=...)} & \inlinecode{DataCollatorWithPadding(tokenizer=...)} \\
\inlinecode{dataset.map(tok,\ batched=True)} & \inlinecode{Trainer(train\_dataset=...)} & \inlinecode{DataLoader(dataset,\ ...)} \\
\inlinecode{DataCollatorWithPadding(...)} & (Trainer가 자동 생성) & \inlinecode{DataLoader(collate\_fn=...)} \\
\inlinecode{with\_format("torch")} & (Trainer가 처리) & \inlinecode{DataLoader} 가 텐서 변환 \\
\bottomrule
\end{tabular}
\end{adjustbox}
\end{booktable}

\textbf{요점}: 이 장에서 익힌 부품들이 9-13장에서 \emph{그대로 입력으로} 들어갑니다. \inlinecode{Trainer} 는 그 부품들을 묶어 학습 루프를 자동화한 것뿐이며, 안에서 일어나는 일은 패턴 B와 같습니다. 커스텀 학습이 필요해지면 패턴 B로 분해해 다시 짤 수 있다는 점이 8장을 손에 익혀두는 가장 큰 이유입니다.



\section{Collator 실습}

DataCollator의 다양한 모습을 직접 돌려봅니다 --- 동적 padding의 효율을 \emph{숫자로} 확인하고, BERT 사전학습에 쓰는 MLM collator를 시연하며, 마지막으로 \emph{직접 작성한 collator} 까지 보면 13장 보조 loss 같은 커스텀 학습이 어떻게 가능한지 감이 잡힙니다.



\subsection{정적 패딩과 동적 패딩}

같은 50 샘플을 두 방식으로 처리해 \emph{전체 토큰 수} 와 \emph{실제 토큰 비율} 을 직접 비교합니다.



\Needspace{26\baselineskip}
\begin{lstlisting}
def tokenize_static(batch):
    return tokenizer(
        batch["text"],
        padding="max_length",
        truncation=True,
        max_length=128,
    )

tokenized_static = small.select(range(50)).map(tokenize_static, batched=True).remove_columns(["text"])
tokenized_static = tokenized_static.with_format(
    "torch",
    columns=["input_ids", "attention_mask", "label"],
)

static_loader = DataLoader(
    tokenized_static,
    batch_size=8,
    shuffle=False,
)

def loader_stats(loader, name):
    real_total = grid_total = 0
    for batch in loader:
        real_total += batch["attention_mask"].sum().item()
        grid_total += batch["attention_mask"].numel()
    return {"method": name, "real_tokens": real_total, "total_tokens": grid_total, "fill_rate": real_total / grid_total}

stats_static = loader_stats(
    static_loader,
    "static (max_length=128)",
)
stats_dyn    = loader_stats(
    dyn_loader,
    "dynamic (DataCollator)",
)

df_stats = pd.DataFrame([stats_static, stats_dyn])
df_stats["fill_rate"] = df_stats["fill_rate"].apply(lambda r: f"{r:.1%}")
print(df_stats.to_string(index=False))

ratio = stats_dyn["total_tokens"] / stats_static["total_tokens"]
print(
    f"\nDynamic vs static total tokens: {ratio:.1%} → ~"
    f"{1-ratio:.0%} reduction"
)
print(
    "(this much self-attention compute saved → faster "
    "training, less memory)"
)
\end{lstlisting}



\subsection{\texorpdfstring{MLM 마스킹}{MLM 마스킹}}

BERT 사전학습은 입력 토큰의 15\%를 \texttt{{[}MASK{]}} 로 가리고 모델이 맞추도록 학습됩니다 (Masked Language Modeling). 그 masking 자체를 담당하는 게 \inlinecode{DataCollatorForLanguageModeling} --- 이 장의 학습엔 안 쓰지만, BERT의 기원을 이해하는 데 직접 보는 게 빠릅니다.



\Needspace{26\baselineskip}
\begin{lstlisting}
from transformers import DataCollatorForLanguageModeling

mlm_collator = DataCollatorForLanguageModeling(
    tokenizer=tokenizer,
    mlm=True,
    mlm_probability=0.15,
)

small_batch = [tokenized_dyn[i] for i in range(3)]
mlm_out = mlm_collator(small_batch)

print(f"input_ids shape: {mlm_out['input_ids'].shape}")
print(
    f"labels shape:    "
    f"{mlm_out['labels'].shape}  (-100 positions ignored by loss)"
)

mask_id = tokenizer.mask_token_id
n_mask = (mlm_out["input_ids"] == mask_id).sum().item()
n_total = mlm_out["input_ids"].numel()
print(
    f"\n[MASK] tokens: {n_mask} / {n_total} ("
    f"{n_mask/n_total:.1%})"
)
print(
    "  (of the 15% masked: 80% [MASK], 10% random token, "
    "10% kept — so [MASK] rate ~12%)"
)
\end{lstlisting}



\Needspace{26\baselineskip}
\begin{lstlisting}
i = 0
orig_ids = small_batch[i]["input_ids"][:25]
mask_ids = mlm_out["input_ids"][i][:25].tolist()
label_ids = mlm_out["labels"][i][:25].tolist()

rows = []
for orig, masked, lbl in zip(orig_ids, mask_ids, label_ids):
    orig_tok = tokenizer.decode([orig])
    masked_tok = tokenizer.decode([masked])
    if lbl == -100:
        label_str = "(ignored)"
    else:
        label_str = tokenizer.decode([lbl])
    changed = "*" if orig != masked else ""
    rows.append({"original": orig_tok, "masked": masked_tok, "label": label_str, "changed": changed})

print(pd.DataFrame(rows).to_string(index=False))
print(
    "\nRows marked * are positions where the collator "
    "masked or replaced the token — the model learns to "
    "predict the original token there."
)
\end{lstlisting}



\subsection{CLM 입력 구성}

같은 \inlinecode{DataCollatorForLanguageModeling} 에 \textbf{\inlinecode{mlm=False}} 만 주면 GPT 같은 \emph{autoregressive 사전학습} 용 collator가 됩니다. 차이는 단순합니다.

\begin{booktable}{8장 CLM 입력 구성 표}{tab:ch08-07}
\begin{adjustbox}{max width=\textwidth}
\begin{tabular}{@{}lll@{}}
\toprule
비교 & MLM (BERT, \inlinecode{mlm=True}) & \textbf{CLM} (GPT, \inlinecode{mlm=False}) \\
\midrule

입력 & 일부 토큰을 \texttt{{[}MASK{]}} 로 가림 & 가리지 않음, 원본 그대로 \\
labels & masked 자리만 원래 토큰, 나머지 \inlinecode{-100} & \inlinecode{input\_ids} 와 동일 (padding은 \inlinecode{-100}) \\
학습 목적 & 양방향 문맥에서 가린 토큰 예측 & 왼쪽 문맥에서 \emph{다음 토큰} 예측 \\
모델 구조 & 양방향 attention & causal mask (왼쪽만 보기) \\
shift 처리 & 없음 & 모델 forward 내부에서 자동 shift-by-one \\
\bottomrule
\end{tabular}
\end{adjustbox}
\end{booktable}

GPT-2 토크나이저는 BERT와 다른 점이 두 가지 있어 약간의 셋업이 필요합니다.

\begin{enumerate}
\def\labelenumi{\arabic{enumi}.}
\tightlist
\item
  \inlinecode{pad\_token} 이 \emph{없습니다}. 사전학습 시 패딩을 안 썼기 때문 --- \inlinecode{eos\_token} 을 pad로 재활용.
\item
  WordPiece 대신 BPE라 \inlinecode{Ġ} (공백) 접두사 표기가 등장 (7장에서 이미 봄).
\end{enumerate}



\Needspace{26\baselineskip}
\begin{lstlisting}
from transformers import AutoTokenizer

gpt_tokenizer = AutoTokenizer.from_pretrained("gpt2")
gpt_tokenizer.pad_token = gpt_tokenizer.eos_token

def gpt_tokenize(batch):
    return gpt_tokenizer(
        batch["text"],
        truncation=True,
        max_length=48,
    )

gpt_tokenized = small.select(range(20)).map(gpt_tokenize, batched=True).remove_columns(["text"])

clm_collator = DataCollatorForLanguageModeling(
    tokenizer=gpt_tokenizer,
    mlm=False,
)

clm_batch = clm_collator([gpt_tokenized[i] for i in range(3)])
print(f"input_ids shape:  {clm_batch['input_ids'].shape}")
print(f"labels shape:     {clm_batch['labels'].shape}")

i = 0
real_len = (clm_batch["attention_mask"][i] == 1).sum().item()
input_first = clm_batch["input_ids"][i][:real_len].tolist()
label_first = clm_batch["labels"][i][:real_len].tolist()
print(
    f"\nFirst sample: input_ids == labels?  (excluding "
    f"padding, all should match)"
)
print(
    f"  Matching positions: {sum(int(a == b) for a, b in zip(input_first, label_first))} / "
    f"{real_len}"
)
\end{lstlisting}



\Needspace{26\baselineskip}
\begin{lstlisting}
i = 0
ids = clm_batch["input_ids"][i][:20].tolist()
lbls = clm_batch["labels"][i][:20].tolist()

rows = []
for tok_id, lbl in zip(ids, lbls):
    tok = gpt_tokenizer.decode([tok_id])
    lbl_str = "(ignored)" if lbl == -100 else gpt_tokenizer.decode([lbl])
    rows.append({"position_token": tok, "label": lbl_str})

print(pd.DataFrame(rows).to_string(index=False))
print(
    "\nObservation: at non-padding positions, input_ids "
    "and labels match."
)
print(
    "    The model shifts by one inside forward() and "
    "predicts the *next token*."
)
print(
    "    e.g. 'Hello world' → at 'Hello' position learn "
    "to predict 'world'."
)
\end{lstlisting}



\subsection{\texorpdfstring{커스텀 collate\_fn}{커스텀 collate\_fn}}

위에서 본 \inlinecode{DataCollatorWithPadding} 와 \inlinecode{DataCollatorForLanguageModeling} 은 \inlinecode{transformers} 가 task별로 미리 만들어준 도구입니다. 그러나 \emph{우리 task} 가 표준 형식에서 벗어나면 (예: 13장 보조 loss처럼 라벨이 두 종류) 직접 함수를 작성해야 합니다.

\inlinecode{collate\_fn} 의 시그니처는 단순합니다 --- 입력은 \emph{샘플 dict의 리스트}, 출력은 \emph{batch dict} (각 키에 stacked 텐서).



\Needspace{26\baselineskip}
\begin{lstlisting}
def custom_collate(batch_list):

    pad_input = collator(
        [{"input_ids": item["input_ids"], "attention_mask": item["attention_mask"]} for item in batch_list]
    )

    raw_labels = torch.tensor(
        [item["label"] for item in batch_list],
        dtype=torch.long,
    )
    aux_scores = raw_labels.float() / 4.0

    pad_input["main_label"] = raw_labels
    pad_input["aux_score"]  = aux_scores
    return pad_input

custom_loader = DataLoader(
    tokenized_dyn,
    batch_size=4,
    shuffle=False,
    collate_fn=custom_collate,
)

batch = next(iter(custom_loader))
print(
    f"input_ids shape:        "
    f"{batch['input_ids'].shape}"
)
print(
    f"attention_mask shape:   "
    f"{batch['attention_mask'].shape}"
)
print(
    f"main_label (int):       "
    f"{batch['main_label'].tolist()}"
)
print(
    f"aux_score (float):      "
    f"{[round(x, 3) for x in batch['aux_score'].tolist()]}"
)
\end{lstlisting}



\textbf{관찰}: \inlinecode{collate\_fn} 한 함수가 \emph{batch 단위 변환의 집결지} 입니다. 라벨 형식 변환, 추가 메타정보 부착, 다중 라벨 dict 등 무엇이든 여기서 처리할 수 있습니다.

13장 보조 loss는 이 패턴을 그대로 사용 --- 메인 라벨(multi-hot)과 보조 라벨(별점 float)을 한 dict 안에 같이 담아서 모델이 두 loss를 합쳐 계산하도록 합니다. \emph{학습 코드의 어디에도 collate\_fn 정의가 없는데 학습이 잘 되네?} 라고 느낀다면, 그건 \inlinecode{Trainer} 가 \inlinecode{tokenizer} 만 보고 \inlinecode{DataCollatorWithPadding} 을 자동 생성한 결과일 가능성이 큽니다 --- 직접 짜야 하는 순간엔 위 패턴이 출발점입니다.



\section{이 장에 등장한 라이브러리·함수}

\subsection{\texorpdfstring{\inlinecode{datasets}}{datasets}}

\begin{booktable}{8장 datasets 표}{tab:ch08-08}
\begin{adjustbox}{max width=\textwidth}
\begin{tabular}{@{}lll@{}}
\toprule
이름 & 한 줄 설명 & 다음 장에서 \\
\midrule

\inlinecode{datasets.load\_dataset} & Hugging Face Hub에서 데이터셋 다운로드 + Apache Arrow 캐시 & 9장 학습 데이터 로드 \\
\inlinecode{Dataset.shuffle(seed)} & 결정론적 셔플 & 재현 가능한 학습 분할 \\
\inlinecode{Dataset.select(indices)} & 지정 인덱스만 선택 & subsample, train/val/test 분할 \\
\inlinecode{Dataset.map(fn,\ batched=True,\ batch\_size=...)} & 배치 단위 변환 (토큰화 등). 결과 자동 캐시 & 9장 입력 전처리 \\
\inlinecode{Dataset.filter(fn)} & 조건 만족 샘플만 & 라벨 필터링, 길이 제한 등 \\
\texttt{Dataset.with\_format(“torch”,\ columns={[}...{]})} & PyTorch tensor 출력 & 모든 학습 장 \\
\bottomrule
\end{tabular}
\end{adjustbox}
\end{booktable}

\subsection{\texorpdfstring{\inlinecode{transformers} 토크나이저 옵션}{transformers 토크나이저 옵션}}

\begin{booktable}{8장 transformers 토크나이저 옵션 표}{tab:ch08-09}
\begin{adjustbox}{max width=\textwidth}
\begin{tabular}{@{}ll@{}}
\toprule
옵션 & 의미 \\
\midrule

\inlinecode{padding=True} & 배치 내 가장 긴 길이까지 padding (동적) \\
\inlinecode{padding="max\_length"} & 항상 max\_length까지 padding (고정) \\
\inlinecode{truncation=True} & max\_length 초과분 잘라냄 \\
\inlinecode{max\_length=N} & 길이 상한 (기본 모델 한도) \\
\inlinecode{return\_tensors="pt"} & PyTorch 텐서로 반환 (\inlinecode{"tf"}, \inlinecode{"np"} 도 가능) \\
\inlinecode{tokenizer.decode(ids)} & ID → 텍스트 역변환 \\
\bottomrule
\end{tabular}
\end{adjustbox}
\end{booktable}

\subsection{PyTorch 데이터 도구}

\begin{booktable}{8장 PyTorch 데이터 도구 표}{tab:ch08-10}
\begin{adjustbox}{max width=\textwidth}
\begin{tabular}{@{}ll@{}}
\toprule
이름 & 한 줄 설명 \\
\midrule

\inlinecode{torch.utils.data.DataLoader} & 배치 + shuffle + multiprocessing \\
\inlinecode{transformers.DataCollatorWithPadding} & DataLoader의 \inlinecode{collate\_fn} --- 배치 시점에 동적 padding \\
\bottomrule
\end{tabular}
\end{adjustbox}
\end{booktable}



\section{체크포인트 질문}

\begin{enumerate}
\def\labelenumi{\arabic{enumi}.}
\tightlist
\item
  \inlinecode{padding=True} 와 \inlinecode{padding="max\_length"} 는 각각 어떤 상황에 적합한가요? 메모리·속도 차이는?
\item
  \inlinecode{attention\_mask=0} 인 자리는 self-attention에서 어떻게 무시되나요? (\inlinecode{-inf} 트릭)
\item
  \inlinecode{datasets} 가 65만 건 데이터를 RAM에 다 안 올리고 인덱싱할 수 있는 이유는 무엇인가요?
\item
  \inlinecode{dataset.map(fn,\ batched=True)} 와 \inlinecode{batched=False} 의 속도 차이는 어디에서 오나요?
\end{enumerate}



\section{FAQ}

\begin{faqBox}{Q1. (실무) \inlinecode{padding=True} 와 \inlinecode{padding="max\_length"} 중 뭘 써야 하나요?}

대부분 \textbf{\inlinecode{padding=True} (동적)} 이 효율적입니다.

\begin{itemize}
\tightlist
\item
  \textbf{\inlinecode{padding=True}}: 배치 안 가장 긴 시퀀스 길이까지만 padding. 짧은 배치는 적은 토큰만 처리해 속도·메모리 절약. 일반 학습/추론에서 표준.
\item
  \textbf{\inlinecode{padding="max\_length"}}: 항상 같은 길이. TPU나 정적 graph(예: TorchScript)처럼 \emph{shape이 매 배치 동일해야 하는} 환경에서. CPU/GPU 학습엔 보통 불필요.
\end{itemize}

\inlinecode{Trainer} 와 \inlinecode{DataCollatorWithPadding} 조합이 동적 padding을 자동 처리하므로, 데이터 전처리 시엔 \textbf{\inlinecode{padding=False}} (또는 생략) + collator에 padding을 맡기는 게 흔한 패턴입니다.

\end{faqBox}

\begin{faqBox}{Q2. (실무) \inlinecode{max\_length} 를 작게 하면 학습이 빨라지지만 성능에 영향은?}

self-attention 비용은 \emph{시퀀스 길이의 제곱} --- \inlinecode{max\_length=128} vs \inlinecode{512} 면 4배 차이가 아니라 \textbf{16배 차이}. 작게 하면 매우 빨라지고 메모리도 4배 절약.

성능 영향은 \textbf{데이터 분포에 따라}:

\Needspace{6\baselineskip}
\begin{lstlisting}
p95 = int(np.percentile(token_lengths, 95))
chosen_max = ((p95 // 32) + 1) * 32
\end{lstlisting}

Yelp는 평균 \textasciitilde150 토큰이라 \inlinecode{max\_length=128} 이 95\% 이상 정상 처리. 더 긴 문서가 많은 데이터(예: 논문 abstract)는 \inlinecode{256} 또는 \inlinecode{512} 가 안전.

\end{faqBox}

\begin{faqBox}{Q3. (이론) \inlinecode{attention\_mask} 가 정확히 모델에서 어떻게 쓰이나요?}

self-attention 계산 직전에 \emph{mask=0인 위치의 점수를 \inlinecode{-inf} 로 바꿉니다}.

\Needspace{8\baselineskip}
\begin{lstlisting}
scores = Q @ K.T / sqrt(d_k)
scores = scores + (1.0 - mask) * -10000
weights = softmax(scores, dim=-1)
output = weights @ V
\end{lstlisting}

softmax 안에서 \texttt{e\^{}(-10000)\ ≈\ 0} 이라 패딩 토큰은 다른 토큰의 표현에 \emph{전혀} 기여하지 않습니다. 학습된 모델 입장에선 패딩이 있든 없든 같은 결과 --- 단지 \emph{계산을 낭비} 한 셈.

\end{faqBox}

\begin{faqBox}{Q4. (실무) 데이터셋이 너무 커서 \inlinecode{map} 이 오래 걸리는데 어떻게 하나요?}

세 가지 가속 기법.

\begin{enumerate}
\def\labelenumi{\arabic{enumi}.}
\tightlist
\item
  \textbf{\inlinecode{batched=True}}: 토크나이저는 batch 호출이 1샘플 호출보다 훨씬 빠름 (Rust 백엔드 활용). 보통 5-10배.
\item
  \textbf{\inlinecode{num\_proc=N}}: 여러 프로세스 병렬화. CPU 코어 수만큼.

\Needspace{9\baselineskip}
\begin{lstlisting}
tokenized = small.map(
    tokenize_fn,
    batched=True,
    num_proc=4,
)
\end{lstlisting}
\item
  \textbf{자동 캐시}: \inlinecode{map} 결과는 디스크에 자동 저장. 같은 함수·같은 데이터로 다시 부르면 즉시 로드 (해시로 식별). Colab 세션이 끝나면 캐시 사라지지만, Drive 마운트로 보존 가능.
\end{enumerate}

\end{faqBox}

\begin{faqBox}{Q5. (이론) \inlinecode{datasets} 가 메모리 효율적인 이유는 무엇인가요?}

핵심: \textbf{Apache Arrow + memory-mapped files}.

\begin{itemize}
\tightlist
\item
  \textbf{Apache Arrow}: column-oriented 바이너리 포맷. 같은 type 데이터를 연속된 메모리에 저장 → cache hit 좋음, 압축 효율 좋음.
\item
  \textbf{memory-mapped (mmap)}: 디스크 파일을 가상 메모리에 \emph{연결만} 해두고 \emph{접근하는 페이지만} 실제 RAM에 올림. 65만 건 데이터셋 객체를 변수에 할당해도 RAM 사용량은 거의 안 늘어남.
\end{itemize}

\Needspace{15\baselineskip}
\begin{lstlisting}
import psutil
process = psutil.Process()
mem_before = process.memory_info().rss / 1024**2

ds = load_dataset("yelp_review_full")

mem_after = process.memory_info().rss / 1024**2
print(
    f"메모리 증가: "
    f"{mem_after - mem_before:.1f} MB  (수십 MB 정도)"
)
\end{lstlisting}

대조적으로 pandas로 같은 데이터를 읽으면 GB 단위 메모리가 필요합니다.

\end{faqBox}

\begin{faqBox}{Q6. (실무) \inlinecode{dataset.set\_format} 과 \inlinecode{with\_format} 의 차이는?}

\begin{itemize}
\tightlist
\item
  \inlinecode{set\_format(type,\ columns)}: \textbf{in-place} 변경. dataset 객체 자체의 출력 형식 변경.
\item
  \inlinecode{with\_format(type,\ columns)}: \textbf{새 dataset 반환}. 원본은 그대로.
\end{itemize}

\Needspace{14\baselineskip}
\begin{lstlisting}
# in-place
tokenized.set_format(
    "torch",
    columns=["input_ids", "attention_mask", "label"],
)

tokenized_torch = tokenized.with_format(
    "torch",
    columns=[...],
)
\end{lstlisting}

대규모 파이프라인에선 \inlinecode{with\_format} 으로 \emph{변환마다 새 객체} 가 안전 (디버깅 시 원본 비교).
\end{faqBox}



\section{오류 실험 (선택)}

다음 코드를 돌려보고, 두 결과의 길이가 왜 다른지 예측해보세요.

\Needspace{16\baselineskip}
\begin{lstlisting}
text = "Hugging Face is amazing!"
out1 = tokenizer(text)
out2 = tokenizer(text, add_special_tokens=False)

print(
    f"기본:                    {len(out1['input_ids'])}, "
    f"{tokenizer.decode(out1['input_ids'])}"
)
print(
    f"add_special_tokens=False: {len(out2['input_ids'])}, "
    f"{tokenizer.decode(out2['input_ids'])}"
)
\end{lstlisting}

힌트: 기본값에선 \texttt{{[}CLS{]}} 와 \texttt{{[}SEP{]}} 가 자동으로 \emph{2개} 추가됩니다 --- \inlinecode{add\_special\_tokens=False} 로 끄면 빠집니다. 분류 작업에선 거의 항상 켜둬야 하지만(BERT 사전학습 시 {[}CLS{]} 자리에 분류 신호가 모이도록 학습됐으므로), 디버깅·연구용으로 끄는 경우가 있습니다.



\begin{previewBox}{미리보기: 다음 장}

\textbf{9장. BERT 회귀 분석 (Regression \& Trainer)}

\begin{itemize}
\tightlist
\item
  Phase 0의 Yelp 별점 회귀(2장)를 \emph{DistilBERT 파인튜닝} 으로 다시 풀기 --- sklearn \inlinecode{LinearRegression} 1초 학습 vs BERT T4 GPU 5-10분 학습
\item
  \textbf{\inlinecode{Trainer} 본격 등장} --- \inlinecode{TrainingArguments}, \inlinecode{compute\_metrics}, evaluation loop
\item
  이 장의 데이터 파이프라인(\inlinecode{datasets.map} + \inlinecode{with\_format("torch")})이 그대로 입력으로 들어감
\item
  Loss는 \inlinecode{MSELoss} (2장과 동일) --- \inlinecode{problem\_type="regression"} 으로 자동 매핑
\item
  \textbf{GPU 필수} --- fp16 옵션 등장
\end{itemize}
\end{previewBox}